% Cut-and-paste TikZ diagram for Section 4.4--4.8
% Requires: \usepackage{tikz} and \usetikzlibrary{arrows.meta,positioning,decorations.pathmorphing,calc,
fit,backgrounds}

\begin{figure}[h!]
\centering
\newcommand{\sep}{7.5}% horizontal distance to right network (adjust as needed)
\begin{tikzpicture}[scale=0.75, every node/.style={transform shape},
  node distance=1.8cm,
  species/.style={circle, draw, thick, minimum size=8mm, inner sep=1pt},
  upf/.style={circle, draw, thick, fill=red!15, minimum size=8mm, inner sep=1pt},
  stab/.style={circle, draw, thick, fill=blue!10, minimum size=8mm, inner sep=1pt},
  rxn/.style={-{Stealth[length=2.5mm]}, thick},
  death/.style={-{Stealth[length=2mm]}, thin, dashed, gray},
  proj/.style={-{Stealth[length=3mm]}, very thick, blue!60!black},
  enlarge/.style={-{Stealth[length=3mm]}, very thick, red!60!black},
  lbl/.style={font=\small, midway},
]

% === LEFT: SIRWS (4 species) ===
\node[stab] (S) at (0, 1.8) {$S$};
\node[upf]  (i) at (1.8, 1.8) {$i$};
\node[upf]  (R) at (1.8, 0) {$R$};
\node[upf]  (W) at (0, 0) {$W$};

% Reactions in SIRWS
\draw[rxn] (S) -- node[lbl, above] {\footnotesize $\beta i S$} (i);
\draw[rxn] (i) -- node[lbl, right] {\footnotesize $\gamma_i i$} (R);
\draw[rxn] (R) -- node[lbl, below] {\footnotesize $\gamma_r R$} (W);
\draw[rxn] (W) -- node[lbl, left]  {\footnotesize $\gamma_w W$} (S);

% Death terms (outgoing dashed)
\draw[death] (S) -- ++(0, 0.7) node[above, font=\tiny] {$\mu_s S$};
\draw[death] (i) -- ++(0.7, 0) node[right, font=\tiny] {$\mu_i i$};
\draw[death] (R) -- ++(0.7, 0) node[right, font=\tiny] {$\mu_r R$};
\draw[death] (W) -- ++(0, -0.7) node[below, font=\tiny] {$\mu_w W$};

% Birth
\draw[rxn, gray] (-0.9, 1.8) node[left, font=\tiny] {$\Lambda$} -- (S);

% Boosting: i+W -> i+R (net effect: W consumed, R produced)
\draw[rxn, dotted] (W) to[bend right=25]
  node[lbl, below, font=\tiny, pos=0.5] {$\nu i W$} (R);

% UPF highlight
\begin{scope}[on background layer]
  \node[draw=red, dashed, rounded corners=6pt,
    fit=(i)(R)(W), inner sep=4pt,
    label={[font=\scriptsize, red!70!black]below left:UPF $\{i,R,W\}$}] {};
\end{scope}

% Label
\node[font=\footnotesize\bfseries, anchor=south] at (0.9, 2.7) {$\mathcal{N}_{\mathrm{SIRWS}}$};

% === RIGHT: Boros--Rost (3 species) ===
\node[stab] (S2) at (\sep, 1.8) {$S\!2$};
\node[upf]  (i2) at (\sep+1.8, 1.8) {$i$};
\node[upf]  (W2) at (\sep+0.9, 0) {$W$};

% Reactions in BR
\draw[rxn] (S2) -- node[lbl, above] {\footnotesize $\beta i S\!2$} (i2);
\draw[rxn] (W2) -- node[lbl, below left, pos=0.4] {\footnotesize $\gamma W$} (S2);
\draw[rxn] (i2) to[bend right=20] node[lbl, right, pos=0.4] {\footnotesize $\nu i W$} (W2);

% Death terms
\draw[death] (S2) -- ++(0, 0.7) node[above, font=\tiny] {$\mu_2 S\!2$};
\draw[death] (i2) -- ++(0.7, 0) node[right, font=\tiny] {$\mu_3 i$};
\draw[death] (W2) -- ++(0, -0.7) node[below, font=\tiny] {$\mu_1 W$};

% Birth
\draw[rxn, gray] (\sep-0.8, 0) node[left, font=\tiny] {$\Lambda$} -- (W2);

% Label
\node[font=\footnotesize\bfseries, anchor=south] at (\sep+0.9, 2.7) {$\mathcal{N}_{\mathrm{BR}}$};

% === ARROWS between networks ===
% Projection (top)
\draw[proj] (2.5, 2.4) -- node[above, font=\small] {$\pi$}
  node[below, font=\scriptsize, text=blue!60!black] {collapse $R$} (\sep-0.7, 2.4);

% Enlargement (bottom)
\draw[enlarge] (\sep-0.7, -0.6) -- node[below, font=\small] {$\mathcal{E}$}
  node[above, font=\scriptsize, text=red!60!black] {inheritance} (2.5, -0.6);

% === Identification annotations ===
\node[font=\scriptsize, text=gray, anchor=west] at (3.8, 1.0) {
  \begin{tabular}{@{}l@{}}
  $S2 \leftrightarrow S$\\
  $W \leftrightarrow W$\\
  $i \leftrightarrow i$\\
  $i{\to}R{\to}W \rightsquigarrow i{\to}W$
  \end{tabular}
};

\end{tikzpicture}

\caption{Contraction and enlargement between the SIRWS network
$\mathcal{N}_{\mathrm{SIRWS}}$ and the Boros--Rost witness
$\mathcal{N}_{\mathrm{BR}}$.
The projection $\pi$ collapses the immune pathway $i\to R\to W$
into the direct interaction $i\to W$, eliminating $R$ and identifying
$S2\leftrightarrow S$.
The red dashed region marks the UPF $\{i,R,W\}$ detected by BiRNe.
The enlargement $\mathcal{E}$ (in the sense of Banaji--Boros--Hofbauer)
reverses this, inserting $R$ as an intermediate and recovering the full
SIRWS dynamics. By the inheritance theorem, Hopf bifurcation transfers
from $\mathcal{N}_{\mathrm{BR}}$ to $\mathcal{N}_{\mathrm{SIRWS}}$.}
\label{fig:contraction}
\end{figure}

% =====================================================================
% Remark on "contractions" in the CRN literature
% (can be placed after the figure or in Section 4.7/4.8)
% =====================================================================

\begin{remark}[Network contractions in the CRN literature]\label{rem:contraction}
The projection $\pi\colon \mathcal{N}_{\mathrm{SIRWS}} \to \mathcal{N}_{\mathrm{BR}}$
that collapses the pathway $i\to R\to W$ into $i\to W$ is an instance of
\emph{species elimination} or \emph{intermediate reduction}, which appears
in several related guises in the CRN literature:

\begin{enumerate}[nosep]
\item \textbf{Intermediate complexes and network translation.}
  Johnston \cite{Johnston14} and Tonello--Johnston \cite{TJ18} develop
  ``network translation" methods that identify intermediate species whose
  elimination preserves qualitative dynamics.
  The collapse of $R$ (which acts as a transient intermediate in the
  immune pathway $i\to R\to W$) fits this framework.

\item \textbf{Variable elimination and quasi-steady-state.}
  Classical QSSA elimination of fast intermediates
  (Briggs--Haldane, Segel--Slemrod \cite{SS89}) contracts pathways
  $A\to B\to C$ to effective reactions $A\to C$ when $B$ equilibrates rapidly.
  Feliu and Wiuf \cite{FW13} formalize this for CRNs, showing that
  intermediate elimination preserves the stoichiometric structure
  relevant for multistationarity.

\item \textbf{Admissible enlargements (the reverse direction).}
  Banaji and Boros \cite{BB22} and Banaji, Boros, and Hofbauer \cite{BBH25}
  define \emph{admissible enlargements}---operations such as species insertion,
  reaction addition, and flow addition that preserve Hopf
  bifurcations via the inheritance theorem.
  Our enlargement $\mathcal{E}\colon \mathcal{N}_{\mathrm{BR}} \to
  \mathcal{N}_{\mathrm{SIRWS}}$ is exactly such a sequence: inserting
  the intermediate $R$, replacing $i+W\to i$ by
  $i+W\to i+R,\; R\to W$, and adding demography.

\item \textbf{Lumping and model reduction.}
  In chemical kinetics, ``lumping" (Wei--Kuo \cite{WK69},
  Li--Rabitz \cite{LR89}) groups species into macro-variables,
  producing a reduced network.
  The identification $\{R,W\}\rightsquigarrow W$ (treating $R$ as a
  transient state of the waning class) is a form of exact lumping
  when the $R\to W$ transition is fast relative to other timescales.
\end{enumerate}

To our knowledge, the term ``contraction" for species-elimination
projections between CRNs has not been standardized.
The closest existing terminology is ``intermediate elimination"
(Feliu--Wiuf) or ``admissible reduction" (viewed as the inverse of
admissible enlargement).
We use ``projection" to emphasize that the map $\pi$ preserves the
local Jacobian feedback structure (the UPF loop) while reducing the
network dimension.
\end{remark}

\begin{remark}[BBH enlargements from $\mathcal{N}_{\mathrm{BR}}$ to
$\mathcal{N}_{\mathrm{SIRWS}}$]\label{rem:enlargements}
The inheritance path $\mathcal{E}\colon \mathcal{N}_{\mathrm{BR}} \to
\mathcal{N}_{\mathrm{SIRWS}}$ decomposes into three admissible enlargements
in the sense of \cite{BBH25}:

\begin{enumerate}[nosep]
\item \textbf{E6 (Splitting reactions).}
  Split the reaction $i+W \to i$ by inserting the new species~$R$
  as an intermediate:
  \[
    i + W \to i \qquad\leadsto\qquad
    i + W \to i + R,\quad R \to W.
  \]
  This introduces one new species ($R$) for one split reaction.
  The stoichiometric submatrix for~$R$ in the two new reactions is
  $(+1,\;-1)$, which has rank~$1$ equal to the number of splits,
  satisfying the E6 rank condition.
  Biologically, this ``unfolds" the collapsed loop
  $i \to W$ in $\mathcal{N}_{\mathrm{BR}}$ into the mechanistic
  pathway $i \to R \to W$ in SIRWS.

\item \textbf{E2 (Fully open extension).}
  Add inflow and outflow reactions for the new species:
  $0 \to R$ and $R \to 0$.
  This ensures the enlarged network is fully open with respect to~$R$,
  providing the dissipation needed for boundedness.

\item \textbf{E1 (New linearly dependent reaction).}
  Add the recovery reaction $i \to R$.
  This is a new reaction among existing species
  (both $i$ and $R$ are present after steps~1--2).
  %It preserves the rank of the stoichiometric matrix
  %since the column $(-1,+1,0,0)^T$ (in the $(S,i,R,W)$ basis)
  %is a linear combination of existing columns.
\end{enumerate}

The E6 step is the structurally decisive one: it is the inverse of the
projection~$\pi$ that collapses $R$ out of the immune pathway.
By the inheritance theorem \cite{BBH25}, each enlargement preserves the
Hopf bifurcation established for $\mathcal{N}_{\mathrm{BR}}$ by
Boros and Rost, and therefore the full SIRWS network inherits Hopf.
\end{remark}

\begin{remark}[Complementary species sets: UPF vs.\ witness]\label{rem:complement}
The BiRNe maximal UPF $\{i,R,W\}$ and the Boros--Rost witness
$\{S,i,W\}$ share the species $\{i,W\}$ but differ in the third:

\begin{center}
\begin{tabular}{lll}
 & \textbf{Species} & \textbf{Role of 3rd species} \\
\hline
BiRNe maximal UPF & $\{i,R,W\}$ & $R$ = intermediate in immune pathway \\
Boros--Rost witness & $\{S,i,W\}$ & $S$ = catalytic environment for infection \\
\end{tabular}
\end{center}

Both are three-species projections of the full Recipe~I oscillatory core,
which consists of the UPF $\{i,R,W\}$ embedded in the stable superset
$\{S,i,R,W\}$.
The two sets arise by dropping complementary species from this
four-species core:
\[
\{S,i,R,W\} \setminus \{S\} = \{i,R,W\} \quad\text{(UPF)},
\qquad
\{S,i,R,W\} \setminus \{R\} = \{S,i,W\} \quad\text{(witness)}.
\]
They play complementary roles:
\begin{itemize}[nosep]
\item The \textbf{UPF} $\{i,R,W\}$ captures the \emph{destabilizing feedback
  loop} $i\to R\to W\to S\to i$: the species whose cyclic interaction
  produces instability at the Jacobian level.
\item The \textbf{witness} $\{S,i,W\}$ captures the \emph{minimal
  mass--action realization}: the species needed for a concrete Hopf
  bifurcation, replacing the intermediate~$R$ with the direct
  interaction $i\to W$ and retaining the susceptible reservoir~$S$
  for autocatalysis ($S+i\to 2i$).
\end{itemize}
The shared pair $\{i,W\}$ is the essential interface where infection meets
waning immunity.
The third species is either the intermediate waystation ($R$, needed for
the Jacobian sign structure) or the susceptible reservoir ($S$, needed for
the mass--action dynamics), depending on whether the analysis is symbolic
(BiRNe) or constructive (Boros--Rost).
\end{remark}

% Suggested references to add to the bibliography:
%
%
%
%
%
%
% [WK69] J. Wei and C.D. Kuo, A lumping analysis in monomolecular
%   reaction systems, Industrial \& Engineering Chemistry Fundamentals 8
%   (1969), 114--123.
%
% [LR89] G. Li and H. Rabitz, A general analysis of exact lumping in
%   chemical kinetics, Chemical Engineering Science 44 (1989), 1413--1430.
